%% The openany option is here just to remove the blank pages before a new chapter
\documentclass[UTF8,a4paper, 10pt, openany]{book}
\usepackage{pagenote}

\usepackage{amsmath,amssymb}
\usepackage{forest}
\usepackage{color}

\usepackage{graphicx}
\graphicspath {{figures/}}
\usepackage{titlepic}
\usepackage{float}

\usepackage{hyperref}
\hypersetup{colorlinks,linkcolor=blue}

\usepackage[margin=1 in]{geometry}
% For customising the endnote markers. Comment these out if you don't want them.
% To prefix each note number with the chapter number
\renewcommand{\thepagenote}{\thechapter-\arabic{pagenote}}

% To have a slightly different formatting for the endnote numbers in the text -- smaller text, sans-serif, square brackets
\renewcommand\notenumintext[1]{\space{\footnotesize\sffamily[FN-#1]}}

% To have a slightly different formatting for the endnote numbers in the notes section. Just the square brackets and sans-serif; normal size.
\renewcommand\notenuminnotes[1]{{\sffamily[FN-#1] }}

% If you want a different name/heading for the end notes
\renewcommand{\notesname}{Eventually, everything is connected.}
%%%%%%%%%%%%% End customisation


%% THIS LINE IS MANDATORY
\makepagenote


% title information
\title{MA Courses Review Notes\\MA2508\\Multi-variable Calculus}
\author{Peilin \textsc{Wu}\\ City University of Hong Kong\\Columbia University School of General Studies}
\titlepic{\includegraphics[scale=0.07]{CityU_logo_2015.png} \includegraphics[scale=0.15]{Columbia_University_School_of_General_Studies_logo.png}}


% the document starts here

\begin{document}
\maketitle
\tableofcontents

\chapter{Vector and the Geometry of Space}
\section{Three-Dimensional Coordinate Systems}
\begin{enumerate}
\item Three-dimensional \textbf{Rectangular} Coordinate System
\\Defined by a triple $(x, y, z)$, can be considered as an orthogonal completed set in three dimension, called the coordinate representation. \\
The gradient expression:
\begin{equation}
 \boxed{\nabla = \mathbf{\hat{i}}\dfrac{\partial }{\partial x}+\mathbf{\hat{j}}\dfrac{\partial }{\partial y}+\mathbf{\hat{k}}\dfrac{\partial }{\partial z}}
\end{equation}

\item Three-dimensional \textbf{Cylindrical} Coordinate System
\\Defined by a triple $(\rho, \varphi, z)$.\\
Relationship between the $(x, y, z)$ system:
\begin{center}
$\begin{array}{cc}
x = \rho \sin \theta  &  \rho ^2 = x^2+y^2\\
y = \rho \cos \theta  &  \tan \theta = \frac{y}{x}
\end{array}$
\end{center}
The gradient expression:
\begin{equation}
 \boxed{\nabla = \mathbf{\hat{r}}\dfrac{\partial }{\partial r}+\dfrac{1}{r}\mathbf{\hat{\theta}}\dfrac{\partial }{\partial \theta}+\mathbf{\hat{z}}\dfrac{\partial }{\partial z}}
\end{equation}

\item Three-dimensional \textbf{Spherical} Coordinate System
\\Defined by a triple $(\rho, \theta, \varphi)$.
\\Relationship between the $(x, y, z)$ system:
\begin{center}
$\begin{array}{cc}
x = \rho \sin \varphi \cos \theta &  \rho ^2 = x^2+y^2+z^2\\
y = \rho \sin \varphi \sin \theta &  \tan \theta = \frac{y}{x}\\
z = \rho \cos \varphi  &  \tan \varphi = \frac{\sqrt{x^2+y^2}}{z}
\end{array}$
\end{center}
The gradient expression:
\begin{equation}
 \boxed{\nabla = \mathbf{\hat{r}}\dfrac{\partial }{\partial r}+\dfrac{1}{r}\mathbf{\hat{\theta}}\dfrac{\partial }{\partial \theta}+\mathbf{\hat{\varphi}}\dfrac{1}{r\sin \theta}\dfrac{\partial }{\partial \varphi}}
\end{equation}
\end{enumerate}

\section{Equation for Line and Plane}
\subsection{Line Function}
Line function given by the vector function is:
\begin{center}
$\mathbf{r}=\mathbf{r_0}+t\mathbf{v}$\\
where $\mathbf{r}=x\mathbf{\hat{i}}+y\mathbf{\hat{j}}+z\mathbf{\hat{k}}$
\end{center}
expressed in the coordinate representation:
\[\begin{cases}
x=x_0+at \\
y=y_0+bt \\
z=z_0+ct
\end{cases}\]
or can be expressed as following:
\begin{center}
$\dfrac{x-x_0}{a}=\dfrac{y-y_0}{b}=\dfrac{z-z_0}{c}=t$
\end{center}
where the vector $\mathbf{v}=a\mathbf{\hat{i}}+b\mathbf{\hat{j}}+c\mathbf{\hat{k}}$ is the directional vector of this line.\\ \\
With this, the angle between any two lines can be expressed as following:\\
Supposing there are two lines $l_1$ and $l_2$:
\begin{center}
line $l_1$: $\dfrac{x-x_0}{a_1}=\dfrac{y-y_0}{b_1}=\dfrac{z-z_0}{c_1}=t_1$\\
line $l_2$: $\dfrac{x-x_0}{a_2}=\dfrac{y-y_0}{b_2}=\dfrac{z-z_0}{c_2}=t_2$
\end{center}
the angle between this two lines equal to the angle between the directional vectors of this two:
\begin{center}
$\cos \theta = \langle \mathbf{v_1},\mathbf{v_2} \rangle =\frac{\mathbf{v_1}\cdot \mathbf{v_2}}{|\mathbf{v_1}||\mathbf{v_2}|}$\\
\end{center}

Therefore, the expression of the angle between two lines are:

\begin{equation}
\boxed{\cos \theta =\dfrac{|a_1a_2+b_1b_2+c_1c_2|}{\sqrt{{a_1}^2+{b_1}^2+{c_1}^2}\sqrt{{a_2}^2+{b_2}^2+{c_1}^2}}}
\end{equation}

\subsection{Plane Function}
Plane function given by the vector function is:
\begin{center}
$\mathbf{n}\cdot \mathbf{r_0}=0$\\
\end{center}
where $\mathbf{r_0}=x_0\mathbf{\hat{i}}+y_0\mathbf{\hat{j}}+z_0\mathbf{\hat{k}}$ ($(x_0,y_0,z_0)$ is a point in the plane) and the $\mathbf{n}=a\mathbf{\hat{i}}+b\mathbf{\hat{j}}+c\mathbf{\hat{k}}$ are the normal vector of this plane
expressed in the coordinate representation:
\begin{center}
$a(x-x_0)+b(y-y_0)+c(z-z_0)=0$
\end{center}
where the vector $\mathbf{n}=a\mathbf{\hat{i}}+b\mathbf{\hat{j}}+c\mathbf{\hat{k}}$ is the normal vector of this plane.\\ \\
With this, the expression of the angle between line and plane is the following:\\
Supposing there are one line $l$ and one plane $n$:
\begin{center}
line $l$: $\dfrac{x-x_0}{a_1}=\dfrac{y-y_0}{b_1}=\dfrac{z-z_0}{c_1}=t$\\
plane $n$: $a_2(x-x_0)+b_2(y-y_0)+c_2(z-z_0)=0$
\end{center}
The angle between them are given by the angle between the directional vector of the line and the normal vector of the plane:
\begin{equation}
\boxed{\sin \varphi =\dfrac{|a_1a_2+b_1b_2+c_1c_2|}{\sqrt{{a_1}^2+{b_1}^2+{c_1}^2}\sqrt{{a_2}^2+{b_2}^2+{c_1}^2}}}
\end{equation}

\subsection{Line function given by two plane}
Notice that the directional vector is perpendicular to two normal vector of two plane $\mathbf{n_1}=a_1\mathbf{\hat{i}}+b_1\mathbf{\hat{j}}+c_1\mathbf{\hat{k}}$ and $\mathbf{n_2}=a_2\mathbf{\hat{i}}+b_2\mathbf{\hat{j}}+c_2\mathbf{\hat{k}}$:
\begin{equation}
\boxed{\mathbf{v}=\mathbf{n_1}\times \mathbf{n_2}=
\begin{vmatrix}
\mathbf{\hat{i}} & \mathbf{\hat{j}} & \mathbf{\hat{k}}\\
a_1 & b_1 & c_1 \\
a_2 & b_2 & c_2
\end{vmatrix}
}
\end{equation}
Moreover the angle between this two plane are:
\begin{equation}
\boxed{\cos \theta =\dfrac{|a_1a_2+b_1b_2+c_1c_2|}{\sqrt{{a_1}^2+{b_1}^2+{c_1}^2}\sqrt{{a_2}^2+{b_2}^2+{c_1}^2}}}
\end{equation}

\subsection{Distance}
Supposing one point $(x_0,y_0,z_0)$ and a plane with its normal vector $\mathbf{n}=a\mathbf{\hat{i}}+b\mathbf{\hat{j}}+c\mathbf{\hat{k}}$ or $ax+by+cz+d=0$:\\
the distance between this two can be expressed as the following:
\begin{equation}
\boxed{d=\dfrac{ax_0+by_0+cz_0+d}{\sqrt{a^2+b^2+c^2}}}
\end{equation}

\section{Cylinders and Quadric Surfaces}

\chapter{Vector Function}
\section{Vector Function}
\subsection{Definition}
Basic form for vector function:
\begin{center}
$\mathbf{v}(t)=\langle f(t),g(t),h(t)\rangle=f(t)\mathbf{\hat{i}}+g(t)\mathbf{\hat{j}}+h(t)\mathbf{\hat{k}}$
\end{center}
\subsection{Limit}
\begin{center}
$\displaystyle	\lim_{t \to t_0} \mathbf{v}(t) = \langle\displaystyle\lim_{t \to t_0}f(t),\displaystyle\lim_{t \to t_0}g(t),\displaystyle\lim_{t \to t_0}h(t)	\rangle$\\
$=\displaystyle\lim_{t \to t_0}f(t)\mathbf{\hat{i}}+\displaystyle\lim_{t \to t_0}g(t)\mathbf{\hat{j}}+\displaystyle\lim_{t \to t_0}h(t)\mathbf{\hat{k}}$
\end{center}
\subsection{Continuity}
$\mathbf{v}(t)$ is continuous at $t=t_0$, iff $\mathbf{v}(t_0)$ is defined.\\
Criteria for the continuity:
\begin{itemize}
\item $\displaystyle\lim_{t \to t_0} \mathbf{v}(t)$ exist
\item $\displaystyle\lim_{t \to t_0} \mathbf{v}(t)=\mathbf{v}(t_0)$
\end{itemize}
\subsection{Derivative}
\begin{center}
$\frac{\mathrm d}{\mathrm d t}\mathbf{v}(t)=\mathbf{v}'(t)=\langle f'(t),g'(t),h'(t) \rangle$\\
$=\frac{\mathrm d}{\mathrm d t}f(t)\mathbf{\hat{i}}+\frac{\mathrm d}{\mathrm d t}g(t)\mathbf{\hat{j}}+\frac{\mathrm d}{\mathrm d t}h(t)\mathbf{\hat{k}}$
\end{center}
\subsection{Operating Rules}
\begin{itemize}
\item $\frac{\mathrm d}{\mathrm d t}(\mathbf{v}(t)\pm \mathbf{w}(t))=\frac{\mathrm d}{\mathrm d t}\mathbf{v}(t)\pm \frac{\mathrm d}{\mathrm d t}\mathbf{w}(t)$
\item $\frac{\mathrm d}{\mathrm d t}[c(t)\mathbf{v}(t)]=[\frac{\mathrm d}{\mathrm d t}c(t)]\mathbf{v}(t)+c(t)\frac{\mathrm d}{\mathrm d t}\mathbf{v}(t)$
\item $\frac{\mathrm d}{\mathrm d t}(\mathbf{v}(t)\cdot \mathbf{w}(t))=\frac{\mathrm d}{\mathrm d t}\mathbf{v}(t)\cdot \mathbf{w}+\mathbf{v}(t)\cdot \frac{\mathrm d}{\mathrm d t}\mathbf{w}(t)$
\item $\frac{\mathrm d}{\mathrm d t}(\mathbf{v}(t)\times \mathbf{w}(t))=\frac{\mathrm d}{\mathrm d t}\mathbf{v}(t)\times \mathbf{w}+\mathbf{v}(t)\times \frac{\mathrm d}{\mathrm d t}\mathbf{w}(t)$
\item $\frac{\mathrm d}{\mathrm d t}\mathbf{v}(c(t))=\mathbf{v}'(c(t))c'(t)$
\end{itemize}
\subsection{Integration}
\begin{center}
$\displaystyle\int\limits_a^b \mathbf{v}(t)dt=\displaystyle\int\limits_a^b f(t)dt \mathbf{\hat{i}}+\displaystyle\int\limits_a^b g(t)dt \mathbf{\hat{j}}+\displaystyle\int\limits_a^b h(t)dt \mathbf{\hat{k}}$
\end{center}
Also similar to the principle of definite integral:\\
If $\mathbf{V}'(t)=\mathbf{v}(t)$, then
\begin{center}
$\displaystyle\int\limits_a^b \mathbf{v}(t)dt=\mathbf{V}(t) |_{a}^b$
\end{center}
\subsection{Parameter Representation a tangent line of a space curve}
\begin{center}
$\overrightarrow{\tau}(t)=\displaystyle\lim_{\Delta t\to 0}\dfrac{\overrightarrow{r}(t+\Delta t)-\overrightarrow{r}(t)}{\Delta t} $
\end{center}
For convenience, we choose to get the unit vector of the tangent vector \\
\begin{center}
$\hat{\tau }(t)=\dfrac{\overrightarrow{\tau }(t)}{|\overrightarrow{\tau }(t)|}$
\end{center}

\subsection{Arc Length and Curvature}
Given a curve $\mathcal{C}$: $x=x(t)\quad y=y(t)\quad z=z(t)$, for small $\Delta t$: the length between two points can be approximated as the distance between two points, thus we have the following formula:
\begin{align*}
\Delta\mathcal{L} & = \sqrt{|\mathbf{r}(t+\Delta t)-\mathbf{r}(t)|^2}\\
& =\sqrt{(x(t+\Delta t)-x(t))^2+(y(t+\Delta t)-y(t))^2+(z(t+\Delta t)-z(t))^2}
\end{align*}

\begin{align*}
\lim_{\Delta t \to 0}\dfrac{\Delta\mathcal{L}}{\Delta t} &= \sqrt{|\mathbf{r}(t+\Delta t)-\mathbf{r}(t)|^2}	\\
&= \sqrt{\dfrac{(x(t+\Delta t)-x(t))^2}{\Delta t^2}+\dfrac{(y(t+\Delta t)-y(t))^2}{\Delta t^2}+\dfrac{(z(t+\Delta t)-z(t))^2}{\Delta t^2}}		\\
&\iff \displaystyle\lim_{\Delta t}\dfrac{\mathcal{L}(t+\Delta t)-\mathcal{L}(t)}{\Delta t}=\sqrt{x'^2(t)+y'^2(t)+z'^2(t)}=\dfrac{d\mathcal{L}}{dt}
\end{align*}

Therefore, we can express the arc length in the differentiation form:

\begin{equation}
\boxed{d\mathcal{L}=\sqrt{x'^2(t)+y'^2(t)+z'^2(t)}dt}
\end{equation}

If we reconsider the tangent vector $\overrightarrow{\tau }(t)$,

\begin{center}
$\overrightarrow{\tau }(t)=x'(t)\mathbf{\hat{i}}+y'(t)\mathbf{\hat{j}}+z'(t)\mathbf{\hat{k}}=\mathbf{r}'(t)$\\
$\dfrac{d\mathcal{L}}{dt}=\sqrt{x'^2(t)\mathbf{\hat{i}}^2+y'^2(t)\mathbf{\hat{j}}^2+z'^2(t)\mathbf{\hat{k}}^2}=|\mathbf{r}'(t)|$
\end{center}

therefore, we have the formula for the arc length:

\begin{equation}
\boxed{\mathcal{L}=s(t)=\displaystyle\int\limits_{t_0}^t|\overrightarrow{\tau }(t)|dt=\displaystyle\int\limits_{t_0}^t|\mathbf{r}'(t)|dt}
\end{equation}

In this way, we can see that $\mathcal{L}=s(t)$ is a function of $t$, suppose that the inverse function exist $t=t(s)$, we can rewrite the curve with respect to $s$:
\begin{center}
$x=x(t)=x(t(s))=X(s)$\\
$y=y(t)=y(t(s))=Y(s)$\\
$z=z(t)=z(t(s))=Z(s)$
\end{center}
as that $s$ is the parameter instead of $t$:
\begin{center}
$\mathcal{L}=s(t)=\displaystyle\int\limits_{t_0}^t|\mathbf{r}'(t(s)|dt \qquad \dfrac{ds}{dt}=|\mathbf{r}'(t)|$\\
$\implies \dfrac{d\mathbf{r}(t(s))}{dt}=\dfrac{d\mathbf{r}}{ds}\dfrac{ds}{dt}=|\mathbf{r}'(t)|\dfrac{d\mathbf{r}}{ds}$\\
$\implies \dfrac{1}{|\mathbf{r}'(t)|}\dfrac{d\mathbf{r}}{dt}=\dfrac{d\mathbf{r}}{ds}$\\
\end{center}
\begin{equation}
\boxed{\implies \hat{\tau}(t)=\dfrac{d\mathbf{r}}{ds}}
\end{equation}
\subsection{Curvature}
Definition:\\
Supposing that $\hat{\tau}$ is the unit tangent vector of the curve $\mathcal{C}$, then the curvature is:
\begin{equation}
\boxed{		\kappa=\left| \dfrac{d\hat{\tau}}{ds} \right|	=	\left|\dfrac{\frac{d\hat{\tau}}{dt}}{\frac{ds}{dt}}\right|	=	\left|\dfrac{\hat{\tau}'(t)}{\mathbf{r}'(t)}\right|	}
\end{equation}

We want to use $\mathbf{r}(t)$ to represent $\hat{\tau}'(t)$, so consider:

\begin{center}
$\mathbf{r}'(t)=|\mathbf{r}'(t)|\hat{\tau}(t)=\dfrac{ds}{dt}\hat{\tau}(t)$\\
$\mathbf{r}''(t)=\dfrac{d^2s}{dt^2}\hat{\tau}(t)+\dfrac{ds}{dt}\hat{\tau}'(t)$\\
$\mathbf{r}'(t)\times \mathbf{r}''(t)=(\dfrac{ds}{dt})^2(\hat{\tau}(t)\times \hat{\tau}'(t))$\\
$|\mathbf{r}'(t)\times \mathbf{r}''(t)|=|(\dfrac{ds}{dt})^2||(\hat{\tau}\times \hat{\tau}'|=|(\dfrac{ds}{dt})^2||\hat{\tau}||\hat{\tau}'|=|(\dfrac{ds}{dt})^2||\hat{\tau}'|$
\end{center}

Thus,

\begin{equation}
\boxed{|\hat{\tau}'|=\dfrac{|\mathbf{r}'(t)\times \mathbf{r}''(t)|}{(\dfrac{ds}{dt})^2}=\dfrac{|\mathbf{r}'(t)\times \mathbf{r}''(t)|}{|\mathbf{r}'(t)|^2}}
\end{equation}
Therefore:
\begin{equation}
\boxed{\kappa=\dfrac{|\hat{\tau}'(t)|}{|\mathbf{r}'(t)|}=\dfrac{|\mathbf{r}'(t)\times \mathbf{r}''(t)|}{|\mathbf{r}'(t)|^3}}
\end{equation}
Note that:\\
for higher dimension general case, this equation always holds true.\\ \\
For $1$-dimensional special case:
\begin{equation}
\boxed{\kappa=\dfrac{|f''(x)|}{[1+(f'(x))^2]^{\frac{3}{2}}}}
\end{equation}
\subsection{The Normal and Binormal Vectors}
The principle unit normal vector:
\begin{center}
$\mathbf{N}(t)=\dfrac{\hat{\tau}'(t)}{|\hat{\tau}'(t)|}$
\end{center}
Osculating plane: the plane that $\hat{\tau}(t)$ and $\mathbf{N}(t)$ span out.\\
Osculating circle: the approximation of the curve by circle with the radius $\rho = \frac{1}{\kappa}$\\
The binormal vector:
\begin{center}
$\mathbf{B}(t)=\hat{\tau}(t)\times \mathbf{N}(t)$
\end{center}
Normal plane: the plane that $\mathbf{B}(t)$ and $\mathbf{N}(t)$ span out.

\chapter{Partial Derivative}
\section{Multi-variables Function}
\subsection{N-dimensional Space}
$n$-th dimensional space can be considered as the $n$-th Cartesian product of $\mathbb{R}$ -- $\mathbb{R}^n$ which can be expressed as the $n$-th triple $(x_1,x_2.\cdots ,x_n)$, and denote the distance between any point in this space and the origin as the norm: $\|\mathbf{x}\|=\sqrt{{x_1}^2+{x_2}^2+\cdots +{x_n}^2}$
\subsection{Multi-variables Function}
A quantity can depend on series variables eq. Temperature $T$ can depend on position $x$ and time $t$:
\begin{center}
$T=T(x,t)$
\end{center}
\textbf{Definition:}\\
Let $\mathcal{D}$ be some region in the $xy$-plane. If there is a rule that assigns to each point in $\mathcal{D}(x,y)$ a unique real number $\mathbb{Z}$, this rule is called a function of there two variable $(x,y)$, denoted as:
\begin{center}
$z=f(x,y)$
\end{center}
\textbf{Associated terminologies:}\\
$\mathcal{D}$ is called the domain of the $f$\\
$\mathcal{R}$ is called the range of the $f$, which is the set of function value\\
\textbf{Note that:} the definition can be immediately extended to the function of more than three variables say:
\begin{center}
$g(x,y,z)$\\
$h(x,y,z,t)$
\end{center}

\subsection{Limit}
\textbf{Definition:}\\
We say that $f(x,y)$ has the limit $\mathcal{L}$ as $(x,y)$ tend to $(x_0,y_0)$ if the following criterion holds:
\\ \\For every $\varepsilon > 0$, there exist a $\delta >0$, such that $|f(x,y)-\mathcal{L}|<\varepsilon $, whenever: $0<\sqrt{(x-x_0)^2+(y-y_0)^2}<\delta$. \\
In this case, we write:
\begin{center}
$\displaystyle\lim_{(x,y) \to (x_0,y_0)} f(x,y) = \mathcal{L}$
\end{center}
If the limit exists, the value will not change along different paths.\\ \\
\textbf{Criteria of the nonexistence} of the limit at one point:\\
If along two different paths, the values corresponding to each path are different, therefore, the limit at the point do not exist.\\
This method can be used to judge existence.\\ \\
\textbf{Example:} Does the $\displaystyle\lim_{(x,y) \to (0,0)} \dfrac{xy^2}{x^2+y^4}$ exist?\\ \\Solution:\\
Along path one $y=x$
\begin{center}
$\displaystyle\lim_{(x,y) \to (0,0)} \dfrac{xy^2}{x^2+y^4}=\displaystyle\lim_{x \to 0} \dfrac{x^3}{x^2+x^4}=0=\mathcal{L}_1$
\end{center}
Along path two $y=-x$
\begin{center}
$\displaystyle\lim_{(x,y) \to (0,0)} \dfrac{xy^2}{x^2+y^4}=\displaystyle\lim_{x \to 0} \dfrac{x^3}{x^2+x^4}=0=\mathcal{L}_2$
\end{center}
Along path three $x=y^2$
\begin{center}
$\displaystyle\lim_{(x,y) \to (0,0)} \dfrac{xy^2}{x^2+y^4}=\displaystyle\lim_{y \to 0}\dfrac{y^4}{2y^4}=\dfrac{1}{2}=\mathcal{L}_3$
\end{center}
Therefore, we see that $\mathcal{L}_1=\mathcal{L}_2 \neq \mathcal{L}_3$, which indicates that the $\lim_{(x,y) \to (0,0)} \dfrac{xy^2}{x^2+y^4}$ does not exist.\\ \\
\textbf{Example:}  Prove that $\displaystyle\lim_{(x,y) \to (0,0)} \dfrac{3x^2y }{x^2+y^2}=0$.\\ \\

Solution: we need to find a $\delta <0$, such that
\begin{center}
$0<\sqrt{x^2+y^2}<\delta $\\
$\left|\dfrac{3x^2y }{x^2+y^2}-0\right|<\varepsilon$
\end{center}
After the scaling to the LHS of the second equation:

\[\left|	\dfrac{3x^2y }{x^2+y^2}	\right| \leq \left|\dfrac{3x^2y^2}{2xy}\right| \leq \dfrac{3}{2}|x| \leq \dfrac{3}{2}\sqrt{x^2+y^2}<\varepsilon \]

Comparing with the first formula, we get $\delta = \frac{2}{3}\varepsilon$\\
Therefore, the limit exist.
\subsection{Continuity}
The $f(x,y)$ is continuous at $(x_0,y_0)$:
\begin{enumerate}
\item $f(x_0,y_0)$ is defined
\item $\displaystyle\lim_{(x,y)\to (x_0,y_0)}f(x,y)=f(x_0,y_0)$
\end{enumerate}
Example: Determine whether\\
\[f(x,y)=\begin{cases}
\dfrac{2xy}{x^2+y^2} & \quad \text{if } (x,y)\neq (0,0)\\
0 & \quad \text{if } (x,y)= (0,0)
\end{cases}\]
is continuous at $(0,0)$ or not?\\ \\
Solution: $f(0,0)=0$ defined
\[\displaystyle\lim_{(x,y)\to (0,0)}\dfrac{2xy}{x^2+y^2}=\begin{cases}
\displaystyle\lim_{x\to 0}\dfrac{2x^2}{2x^2}=\dfrac{2}{3}=\mathcal{L}_1 & \quad \text{path 1 } x=y\\
\displaystyle\lim_{x\to 0}\dfrac{2x(-x)}{2x^2}=-\dfrac{2}{3}=\mathcal{L}_2 & \quad \text{path 2 } -x=y\\
\end{cases}\]
$\mathcal{L}_1\neq \mathcal{L}_2$ indicates that the limits does not exist, therefore, the function is discontinuous at this point.\\
Note: Continuity and limit can be expanded into higher dimension.

\section{Definition}
Partial derivative defined as following:
\begin{center}
$\dfrac{\partial f}{\partial x}=\displaystyle\lim_{\Delta x \to 0}\dfrac{f(x+\Delta x,y)-f(x,y)}{\Delta x}$
\end{center}
denoted as: $f_x(x,y)$ or $\dfrac{\partial f(x,y)}{\partial x}$ or $z_x$.
\subsection{Calculation of the partial derivatives}
To calculate $\dfrac{\partial f}{\partial x}$, we regard the $y$ as a constant and then take the derivative w.r.t $x$.
\subsection{Partial Derivatives Family}
For the function of $z=f(x,y)$
\begin{center}
$\dfrac{\partial f}{\partial x}=\displaystyle\lim_{\Delta x \to 0}\dfrac{f(x+\Delta x,y)-f(x,y)}{\Delta x}$\\
$\dfrac{\partial f}{\partial y}=\displaystyle\lim_{\Delta y \to 0}\dfrac{f(x,y+\Delta y)-f(x,y)}{\Delta y}$
\end{center}

\subsection{Example}
Find the two partial derivatives of the function: $f(x,y)=\sqrt{2xy^2+3y}$ and the find their value at $(3,3)$\\
Solution:\\

\begin{center}
$\dfrac{\partial f}{\partial x}=\dfrac{1}{2}\dfrac{\dfrac{\partial}{\partial x}(2xy^2+3y)}{\sqrt{2xy^2+3y}}=\dfrac{2y^2}{2\sqrt{2xy^2+3y}}$\\
$\dfrac{\partial f}{\partial y}=\dfrac{1}{2}\dfrac{\dfrac{\partial}{\partial y}(2xy^2+3y)}{\sqrt{2xy^2+3y}}=\dfrac{4xy+3}{2\sqrt{2xy^2+3y}}$\\
$\displaystyle\dfrac{\partial f}{\partial x}|_{(3,3)}=\dfrac{3}{\sqrt{7}}\qquad \displaystyle\dfrac{\partial f}{\partial y}|_{(3,3)}=\dfrac{13}{2\sqrt{7}}$
\end{center}

\section{Higher Partial Derivatives}
Partial derivatives can be easily expanded into higher order and higher dimensions.\\ \\
For $g(x,y,z)$, we can compute its partial derivatives as the following formula:
\begin{center}
$\dfrac{\partial g(x,y,z)}{\partial x}=\displaystyle\lim_{\Delta x\to 0} \dfrac{g(x+\Delta x,y,z)-g(x,y,z)}{\Delta x}$\\
$\dfrac{\partial g(x,y,z)}{\partial y}=\displaystyle\lim_{\Delta y\to 0} \dfrac{g(x,y+\Delta y,z)-g(x,y,z)}{\Delta y}$\\
$\dfrac{\partial g(x,y,z)}{\partial z}=\displaystyle\lim_{\Delta z\to 0} \dfrac{g(x,y,z+\Delta z)-g(x,y,z)}{\Delta z}$
\end{center}
For given $z=f(x,y)$, higher order of the partial derivatives are given as the following:
\begin{center}
$\dfrac{\partial z}{\partial x}=\dfrac{\partial f(x,y)}{\partial x}=g(x,y)$\\
$\dfrac{\partial z}{\partial y}=\dfrac{\partial f(x,y)}{\partial y}=h(x,y)$
\end{center}
Therefore, for the second partial derivative:
\begin{center}
$\dfrac{\partial g(x,y)}{\partial x}=\dfrac{\partial}{\partial x}\dfrac{\partial f(x,y)}{\partial x}=\dfrac{\partial^2 f(x,y)}{\partial x^2}=f_{xx}(x,y)$\\
$\dfrac{\partial h(x,y)}{\partial y}=\dfrac{\partial}{\partial y}\dfrac{\partial f(x,y)}{\partial y}=\dfrac{\partial^2 f(x,y)}{\partial y^2}=f_{yy}(x,y)$\\
$\dfrac{\partial g(x,y)}{\partial y}=\dfrac{\partial}{\partial y}\dfrac{\partial f(x,y)}{\partial x}=\dfrac{\partial^2 f(x,y)}{\partial x\partial y}=f_{xy}(x,y)$\\
$\dfrac{\partial h(x,y)}{\partial x}=\dfrac{\partial}{\partial x}\dfrac{\partial f(x,y)}{\partial y}=\dfrac{\partial^2 f(x,y)}{\partial y\partial x}=f_{yx}(x,y)$
\end{center}
About the last two second order partial derivatives symmetric to each other, their relationship can be determined by the Clairant's Theorem.

\section{Clairant's Theorem}
If $f(x,y)$ is defined in a region that contain $(x_0,y_0)$, then given that $f_{xy}$ and $f_{yx}$ are continuous at $(x_0,y_0)$, by this, the following formula hold true to itself:
\begin{equation}
\boxed{f_{xy}(x_0,y_0)=f_{yx}(x_0,y_0)}
\end{equation}
Noted that: In MA2508 and MA2158, if not states, we always assume the condition in the theorem are satisfied: $f_{xy}=f_{yx}$.\\ \\
In general, by Clairant's Theorem, higher order partial derivatives can be defined as following:
\begin{center}
$\dfrac{\partial^{m+n}f(x,y)}{\partial x^m \partial y^n}$ or $\dfrac{\partial^{a+b+c}g(x,y,z)}{\partial x^a \partial y^b \partial z^c}$ or even higher.
\end{center}
Noted: For higher derivatives, computing order is important.\\ \\
Example: Given
\begin{center}
$f(x,y)=(1+xy)\ln (1+x^2)$
\end{center}
Find
\begin{center}
$\dfrac{\partial^{8}f(x,y)}{\partial x^5 \partial y^3}$
\end{center}
Solution:
\[
  \begin{cases}
    \dfrac{\partial^{8}f(x,y)}{\partial x^5 \partial y^3}=\dfrac{\partial^{7}f(x,y)}{\partial x^4 \partial y^3}\dfrac{\partial f(x,y)}{\partial x}=\dfrac{\partial^{7}f(x,y)}{\partial x^4 \partial y^3}(y\ln (1+x^3)+(1+xy)\dfrac{2y}{1+x^2})=\cdots \\
    \dfrac{\partial^{8}f(x,y)}{\partial x^5 \partial y^3}=\dfrac{\partial^{7}f(x,y)}{\partial x^5 \partial y^2}\dfrac{\partial f(x,y)}{\partial y}=\dfrac{\partial^{7}f(x,y)}{\partial x^5 \partial y^2}(x\ln (1+x^2))=0 \\
  \end{cases}\]

\section{Geometrical Representation}
\subsection{Tangent Line}
The tangent line is given by the tangent vector on the surface $z=f(x,y)$ at the point $(x_0,y_0)$. However, there are infinite number of tangent line at one point that consist a tangent plane. So firstly, let's consider the tangent vector along the fixed $y$ section curve of the surface, the slope of this curve is given by the following:
\begin{center}
$k_y=\dfrac{\partial f}{\partial x}|_{(x_0,y_0)}$\\
$k_x=\dfrac{\partial f}{\partial y}|_{(x_0,y_0)}$
\end{center}
Therefore, we can express the two tangent line of the surface as following;
\begin{center}
$l_{y}: \quad z-z_0=\dfrac{\partial f}{\partial x}|_{(x_0,y_0)}(x-x_0)$\\
$l_{x}: \quad z-z_0=\dfrac{\partial f}{\partial y}|_{(x_0,y_0)}(y-y_0)$
\end{center}
Furthermore, if we know the two directional tangent vector, we can easily compute the tangent plane.
\subsection{Tangent Plane}
Given the two vectors of the above two directional vector, we can express the tangent plane as following: \\For surface $z=f(x,y)$
\begin{equation}
\boxed{z-f(x_0,y_0)=\dfrac{\partial f}{\partial x}|_{(x_0,y_0)}(x-x_0)+\dfrac{\partial f}{\partial y}|_{(x_0,y_0)}(y-y_0)}
\end{equation}
For the surface with implicit function $F(x,y,z)=0$, the tangent plane:
\begin{equation}
\boxed{\dfrac{\partial F}{\partial x}|_{(x_0,y_0,z_0)}(x-x_0)+\dfrac{\partial F}{\partial y}|_{(x_0,y_0,z_0)}(y-y_0)\dfrac{\partial F}{\partial z}|_{(x_0,y_0,z_0)}(z-z_0)=0}
\end{equation}
Example: Find the equation for the tangent plane of the surface $z=\sin (x+y)$ at the point $(1,-1,0)$
\\Solution:
\begin{center}
$\left. \dfrac{\partial z}{\partial x} \right|_{(1,-1,0)}=1 \qquad \left. \dfrac{\partial z}{\partial y} \right|_{(1,-1,0)}=1$\\
$z=(x-1)+(y+1)=x+y$
\end{center}

\section{Partial differential}
\subsection{Recall}
If $y=f(x)$ is differentiable, that is $f(a+dx)=f(a)+\Delta y=f(a)+f'(a)dx+\varepsilon dx$, where $dx$ is the increment with $\varepsilon \to 0$ as $dx \to 0$, then $dy|_{a}=f'(a)dx$ is called differential.\\ \\
We can see that:
\begin{center}
$f(a+dx)-f(a)\approx dy|_{a}$ for small $dx$\\
$f(a+dx)\approx f(a)+dy|_{a}$
\end{center}
The value at $a+dx$ can be approximately determined.
\subsection{Definition}
$z=f(x,y)$ is said to be differentiable at $(x_0,y_0)$ if:
\begin{center}
$dz|_{(x_0,y_0)}=f(x+dx,y+dy)-f(x,y)=f_x(x_0,y_0)dx+f_y(x_0,y_0)+\varepsilon_1dx+\varepsilon_2dy$
\end{center}
where $(\varepsilon_1,\varepsilon_2)\to (0,0)$ when $(dx,dy)\to (0,0)$.\\
We call $dz|_{(x_0,y_0)}=f(x+dx,y+dy)-f(x,y)=f_x(x_0,y_0)dx+f_y(x_0,y_0)$ to be the differential of $f$ at $(x_0,y_0)$
\subsection{Linear Approximation}
\begin{center}
$f(x+dx,y+dy)\approx f(x_0,y_0)+f_x(x_0,y_0)dx+f_y(x_0,y_0)=f(x_0,y_0)+dz|_{(x_0,y_0)}$\\
$dz|_{(x_0,y_0)}=\left. \dfrac{\partial f}{\partial x} \right|_{(x_0,y_0)}dx + 	\left.\dfrac{\partial f}{\partial y} \right|_{(x_0,y_0)}dy$
\end{center}
\subsection{Criteria for Differential}
If $f_x$ and $f_y$ exist in a neighborhood of $(x_0,y_0)$ and continuous at $(x_0,y_0)$, then $f(x,y)$ is differentiable at $(x_0,y_0)$.

\section{The Chain Rule}
\subsection{Recall and Analogy}
$y=f(x)$, $x=g(t)$, the derivative of $y$ w.r.t $t$ is following:
\begin{center}
$\dfrac{dy}{dt}=\dfrac{dy}{dx}\dfrac{dx}{dt}$
\end{center}
Therefore, for differentiable function $z=f(x,y)$, where $x=g(t)$ and $y=h(t)$, we have the following formula:
\begin{center}
$\dfrac{dz}{dt}=\dfrac{\partial f}{\partial x}\dfrac{dx}{dt}+\dfrac{\partial f}{\partial y}\dfrac{dy}{dt}$
\end{center}
\subsection{Proof}
An increment $\Delta t$ leads to the increment $\Delta x$, $\Delta y$ and $\Delta z$.
\begin{center}
$\dfrac{dz}{dt}=\displaystyle\lim_{\Delta t\to 0}\dfrac{\Delta z}{\Delta t}=\displaystyle\lim_{\Delta t\to 0}\dfrac{1}{\Delta t}(f_x(x,y)dx+f_y(x,y)+\varepsilon_1dx+\varepsilon_2dy)$\\
$=\displaystyle\lim_{\Delta t\to 0}\dfrac{\Delta x}{\Delta t}f_x(x,y)+\displaystyle\lim_{\Delta t\to 0}\dfrac{\Delta y}{\Delta t}f_y(x,y)+\displaystyle\lim_{\Delta t\to 0}\varepsilon_1\lim_{\Delta t\to 0}\dfrac{\Delta x}{\Delta t}+\displaystyle\lim_{\Delta t\to 0}\varepsilon_2\lim_{\Delta t\to 0}\dfrac{\Delta y}{\Delta t}$\\
$=f_x(x,y)\dfrac{dx}{dt}+f_y(x,y)\dfrac{dy}{dt}=\dfrac{\partial f}{\partial x}\dfrac{dx}{dt}+\dfrac{\partial f}{\partial y}\dfrac{dy}{dt}$
\end{center}
\subsection{Tree diagram}
Procedure for writing down the chain formula:\\
To calculate the derivative of $z$ w.r.t $t$, we can draw the tree diagram and count how many different path from $z$ to $t$. For each path, calculate the product of derivative w.r.t the local variable. Finally, add all the product together.

\begin{center}
\begin{forest}
  for tree={
    fit=band,% spaces the tree out a little to avoid collisions
  }
  [z[x[t]][y[t]]]
\end{forest}\\
$\dfrac{dz}{dt}=\dfrac{\partial f}{\partial x}\dfrac{dx}{dt}+\dfrac{\partial f}{\partial y}\dfrac{dy}{dt}$
\end{center}

Example: Given $xy+yz+zx=u$, where $x=st$, $y=\exp{st}$ and $z=t^2$, find $\dfrac{dy}{ds}$\\ \\
Solution:
\begin{center}
\begin{forest}
  for tree={
    fit=band,% spaces the tree out a little to avoid collisions
  }
  [u[x[s][t]][y[s][t]][z[t]]]
\end{forest}\\
$\dfrac{du}{ds}=\dfrac{\partial u}{\partial x}\dfrac{dx}{ds}+\dfrac{\partial u}{\partial y}\dfrac{dy}{ds}=(y+z)t+(x+z)t\exp{st}$
\end{center}
Example: Given: $z=\dfrac{x}{y}$, $x=r\exp{st}$, $y=rs\exp{t}$, find $\dfrac{\partial^2 z}{\partial t\partial r}$\\ \\
Solution:
\begin{center}
\begin{forest}
  for tree={
    fit=band,% spaces the tree out a little to avoid collisions
  }
  [z[x[r][s][t]][y[r][s][t]]]
\end{forest}\\
$\dfrac{\partial z}{\partial t}=\dfrac{\partial z}{\partial x}\dfrac{\partial x}{\partial t}+\dfrac{\partial z}{\partial y}\dfrac{\partial y}{\partial t}=\dfrac{1}{y}rs\exp{st}+(-\dfrac{x}{y^2}rs\exp{st})$
\end{center}
Then draw the second diagram:
\begin{center}
\begin{forest}
  for tree={
    fit=band,% spaces the tree out a little to avoid collisions
  }
  [$\dfrac{\partial z}{\partial t}$ [x[r][s][t]][y[r][s][t]][r][s][t]]
\end{forest}\\
$\dfrac{\partial^2 z}{\partial t\partial r}=\dfrac{\partial }{\partial r}\dfrac{\partial z}{\partial t}=\dfrac{\partial }{\partial x}(\dfrac{\partial z}{\partial t})\dfrac{\partial x}{\partial r}+\dfrac{\partial }{\partial y}(\dfrac{\partial z}{\partial t})\dfrac{\partial y}{\partial r}+\dfrac{\partial }{\partial r}(\dfrac{\partial z}{\partial t})=\cdots $
\end{center}
\textbf{Example: (Important)} If $z=f(x,y)$, where $x=r\cos \theta $, $y=r\sin \theta $ show that:
\[\left( \dfrac{\partial z}{\partial x} \right)^2+(\dfrac{\partial z}{\partial x})^2=\left(\dfrac{\partial z}{\partial r}\right)^2+\dfrac{1}{r^2}\left( \dfrac{\partial z}{\partial \theta }\right)^2\]

Solution:
\begin{center}
\begin{forest}
  for tree={
    fit=band,% spaces the tree out a little to avoid collisions
  }
  [z[x[r][$\theta $]][y[r][$\theta $]]]
\end{forest}
\end{center}
Path 1: from the R.H.S to the L.H.S
\begin{center}
$\dfrac{\partial z}{\partial r}=\dfrac{\partial z}{\partial x}\dfrac{\partial x}{\partial r}+\dfrac{\partial z}{\partial y}\dfrac{\partial y}{\partial r}=\cos\theta \dfrac{\partial z}{\partial x}+\sin\theta \dfrac{\partial z}{\partial y}$ (Eq1)\\
$\dfrac{\partial z}{\partial \theta}=\dfrac{\partial z}{\partial x}\dfrac{\partial x}{\partial \theta}+\dfrac{\partial z}{\partial y}\dfrac{\partial y}{\partial \theta}=-r\sin\theta \dfrac{\partial z}{\partial x}+r\cos\theta \dfrac{\partial z}{\partial y}$ (Eq2)\\
$R.H.S=(\dfrac{\partial z}{\partial r})^2+\dfrac{1}{r^2}(\dfrac{\partial z}{\partial \theta })^2=(\cos\theta \dfrac{\partial z}{\partial x}+\sin\theta \dfrac{\partial z}{\partial y})^2+\dfrac{1}{r^2}(-r\sin\theta \dfrac{\partial z}{\partial x}+r\cos\theta \dfrac{\partial z}{\partial y})^2=(\dfrac{\partial z}{\partial x})^2\cos^2\theta+2\dfrac{\partial z}{\partial y}\cos\theta \sin\theta +(\dfrac{\partial z}{\partial x})^2\sin^2\theta +(\dfrac{\partial z}{\partial x})^2\sin^2\theta-2\dfrac{\partial z}{\partial y}\cos\theta \sin\theta +(\dfrac{\partial z}{\partial x})^2\cos^2\theta =L.H.S$
\end{center}
Path 2: from the L.H.S to the R.H.S
Eq1 $\times r\cos \theta -$ Eq2 $\times r\sin \theta$\\
Eq1 $\times r\sin \theta +$ Eq2 $\times r\cos \theta$\\
This yields:
\begin{center}
$\dfrac{\partial z}{\partial x}=\dfrac{\partial z}{\partial r}\cos\theta -\dfrac{\partial z}{\partial \theta}\dfrac{\sin\theta}{r}$\\
$\dfrac{\partial z}{\partial y}=\dfrac{\partial z}{\partial r}\sin\theta +\dfrac{\partial z}{\partial \theta}\dfrac{\cos\theta}{r}$\\
$L.H.S=(\dfrac{\partial z}{\partial x})^2+(\dfrac{\partial z}{\partial x})^2=\cdots=R.H.S$
\end{center}

\section{Implicit Differentiation}
Given: $F(x,y)=0$, and say $y$ can be regarded as an implicit function of $x$. Find the $\dfrac{dy}{dx}$
\begin{center}
\begin{forest}
  for tree={
    fit=band,% spaces the tree out a little to avoid collisions
  }
  [$F$[x][y[x]]]
\end{forest}
\end{center}
Differentiating w.r.t $x$:
\begin{center}
$\dfrac{\partial F}{\partial x}+\dfrac{\partial F}{\partial y}\dfrac{dy}{dx}=0$\\
$\implies \dfrac{dy}{dx}=-\dfrac{\dfrac{\partial F}{\partial x}}{\dfrac{\partial F}{\partial y}}$
\end{center}
Given: $F(x,y,z)=0$, and say $z$ can be regarded as an implicit function of $x$ and $y$. Find the $\dfrac{\partial z}{\partial x}$ and $\dfrac{\partial z}{\partial x}$
\begin{center}
\begin{forest}
  for tree={
    fit=band,% spaces the tree out a little to avoid collisions
  }
  [$F$[x][y][z[x][y]]]
\end{forest}
\end{center}
Differentiating w.r.t $x$:
\begin{center}
$\dfrac{\partial F}{\partial x}+\dfrac{\partial F}{\partial z}\dfrac{\partial z}{\partial x}=0$\\
$\implies \dfrac{\partial z}{\partial x}=-\dfrac{\dfrac{\partial F}{\partial x}}{\dfrac{\partial F}{\partial z}}$
\end{center}
Similarly, Differentiating w.r.t $y$:
\begin{center}
$\dfrac{\partial F}{\partial y}+\dfrac{\partial F}{\partial z}\dfrac{\partial z}{\partial y}=0$\\
$\implies \dfrac{\partial z}{\partial y}=-\dfrac{\dfrac{\partial F}{\partial y}}{\dfrac{\partial F}{\partial z}}$
\end{center}

\section{Directional Derivatives and Gradient Vector}
Consider $z=f(x,y)$:
\begin{center}
$f_x(x_0,y_0)=\displaystyle\lim_{h\to 0}\dfrac{f(x_0+h,y_0)-f(x_0,y_0)}{h}$\\ (net change along the $x$-direction)\\
$f_x(x_0,y_0)=\displaystyle\lim_{h\to 0}\dfrac{f(x_0,y_0+h)-f(x_0,y_0)}{h}$\\ (net change along the $y$-direction)
\end{center}
Therefore, naturally, what is the net change in any arbitrary direction $\mathbf{\hat{u}}=a\mathbf{\hat{i}}+b\mathbf{\hat{j}}$
\begin{center}
$x_1=x_0+h\cos\theta \qquad a=\mathbf{\hat{u}}\cdot \mathbf{\hat{i}}$\\
$y_1=y_0+h\sin\theta \qquad b=\mathbf{\hat{u}}\cdot \mathbf{\hat{j}}$\\
$f_u(x_0,y_0)=\displaystyle\lim_{h\to 0}\dfrac{f(x_0+ah,y_0+bh)-f(x_0,y_0)}{h}$
\end{center}
If the limit exists, we call it to be the directional derivatives of this function $f(x,y)$ at $(x_0,y_0)$ along $\mathbf{\hat{u}}$-direction and denoted it by $D_{\mathbf{\hat{u}}}f(x_0,y_0)$

\subsection{Theorem}
If $f$ is differentiable function of $x$ and $y$
\begin{center}
$D_{\mathbf{\hat{u}}}f(x_0,y_0)=f_x(x_0,y_0)a+f_y(x_0,y_0)b \qquad \mathbf{\hat{u}}=a\mathbf{\hat{i}}+b\mathbf{\hat{j}}$
\end{center}
Proof: define $f(x_0+ha,y_0+hb)=g(h)$\\
Then $g(0)=f(x_0,y_0)$, by definition:
\begin{center}
$D_{\mathbf{\hat{u}}}f(x_0,y_0)=\displaystyle\lim_{h\to 0}\dfrac{g(h)-g(0)}{h}=\left.\dfrac{dg}{dh}\right|_{h=0}$
\end{center}
Let $x=x_0+ah$ and $y=y_0+bh$ and differentiating w.r.t $h$:
\begin{center}
\begin{forest}
  for tree={
    fit=band,% spaces the tree out a little to avoid collisions
  }
  [$g$[x[h]][y[h]]]
\end{forest}
\end{center}
\begin{center}
$\dfrac{dg}{dh}|_{h=0}=\dfrac{\partial f}{\partial x}|_{(x_0,y_0)}\dfrac{dx}{dh}|_{h=0}+\dfrac{\partial f}{\partial y}|_{(x_0,y_0)}\dfrac{dy}{dh}|_{h=0}$\\
$=f_x(x_0,y_0)a+f_y(x_0,y_0)b$
\end{center}
\textbf{Observation:}\\
We can write the directional derivatives as the dot product of the gradient vector and the directional vector:
\begin{center}
$D_{\mathbf{\hat{u}}}f(x_0,y_0)=\nabla f(x_0,y_0)\cdot \mathbf{\hat{u}}$\\
$\nabla f(x_0,y_0)=\dfrac{\partial f}{\partial x}|_{(x_0,y_0)}\mathbf{\hat{i}}+\dfrac{\partial f}{\partial y}|_{(x_0,y_0)}\mathbf{\hat{j}}$\\
Where the operator: $\nabla = \dfrac{\partial }{\partial x}\mathbf{\hat{i}}+\dfrac{\partial }{\partial y}\mathbf{\hat{j}}$
\end{center}
Note: This gradient can be extended to higher dimension.
\subsection{Example}
Find the directional derivative of $g(x,y)=\exp{x}\cos y$ at $(1,\dfrac{\pi}{6})$ in the directional vector $\mathbf{v}=\mathbf{\hat{i}}-\mathbf{\hat{j}}$:\\
Solution:
\begin{center}
$\mathbf{\hat{v}}=\dfrac{\mathbf{v}}{|\mathbf{v}|}=\dfrac{\mathbf{\hat{i}}-\mathbf{\hat{j}}}{\sqrt{1^2+1^2}}=\dfrac{\mathbf{\hat{i}}}{\sqrt{2}}-\dfrac{\mathbf{\hat{j}}}{\sqrt{2}}$\\
$\nabla g(x,y)=\dfrac{\partial g}{\partial x}\mathbf{\hat{i}}+\dfrac{\partial g}{\partial y}\mathbf{\hat{j}}=\exp{x}(\cos y\mathbf{\hat{i}}-\sin y\mathbf{\hat{j}})$\\
$\nabla g(1,\dfrac{\pi}{6})=e (\dfrac{\sqrt{3}\mathbf{\hat{i}}}{2}-\dfrac{\mathbf{\hat{j}}}{2})$
\end{center}
Then, after the dot product
\begin{center}
$\nabla g(1,\dfrac{\pi}{6})\cdot \mathbf{\hat{v}}=\dfrac{\sqrt{3}+1}{2\sqrt{2}}e$
\end{center}
\subsection{Extension}
To function of three variables (dimension)
\begin{center}
$\hat{u}=a\hat{i}+b\hat{j}+c\hat{k} \qquad w=f(x,y,z)$\\
$\implies D_{\hat{u}}f(x,y,z)=\displaystyle\lim_{h\rightarrow 0}=\dfrac{f(x+ha,y+hb,z+,hc)-f(x,y,z)}{h}$\\
$\nabla f(x,y,z)=\dfrac{\partial f}{\partial x}\hat{i}+\dfrac{\partial f}{\partial y}\hat{j}+\dfrac{\partial f}{\partial z}\hat{k}$\\
$\boxed{D_{\hat{u}}f(x,y,z)=\nabla f(x,y,z)\cdot \hat{u}}$
\end{center}
\subsection{Fastest Changing Direction}
\textbf{Theorem: }The maximum value of the $D_{\hat{u}}f(\mathbf{v})$ is the $|\nabla f(\mathbf{r})|$\\ \\
\textbf{Proof: }\\
\begin{center}
$D_{\hat{u}}f(\mathbf{r})=\nabla f(\mathbf{r})\cdot \hat{u}$\\
$= |\nabla f(\mathbf{r})|\cdot |\hat{u}|\cos \theta \leq |\nabla f(\mathbf{r})|$
\end{center}
This suggests: along the gradient direction, the function changes fastest.
\subsection{Geometrical Meaning of the Gradient Vector}
Consider a surface $\mathcal{S}$ given by:
\begin{center}
$F(x,y,z)=0$
\end{center}
Then, we consider an arbitrary curve $\mathcal{C}$ on $\mathcal{S}$ surface that suppose that $\mathcal{C}$: $x=x(t),y=y(t),z=z(t)$
\begin{center}
$F(x(t),y(t),z(t))=0$
\end{center}
Differentiating this equation w.r.t t
\begin{center}
$\dfrac{\partial F}{\partial x}\dfrac{dx}{dt}+\dfrac{\partial F}{\partial y}\dfrac{dy}{dt}+\dfrac{\partial F}{\partial z}\dfrac{dz}{dt}=0$
\end{center}
Then, we denote:
\begin{align*}
\mathbf{r}'(t)&=\dfrac{dx}{dt}\mathbf{\hat{i}}+\dfrac{dy}{dt}\mathbf{\hat{j}}+\dfrac{dz}{dt}\mathbf{\hat{k}}\\
\nabla F(x,y,z)&=\dfrac{\partial F}{\partial x}\mathbf{\hat{i}}+\dfrac{\partial F}{\partial y}\mathbf{\hat{j}}+\dfrac{\partial F}{\partial z}\mathbf{\hat{k}}\\
\end{align*}
After this, rewrite the differentiating equation:
\begin{equation}
\boxed{
\nabla F(x,y,z)\cdot \mathbf{r}'(t)=0
}
\end{equation}
This means that: $\nabla F(x,y,z)$ is the normal vector of the tangent plane since $\nabla F(x,y,z)$ is perpendicular to any tangent vector along an arbitrary curve.\\ \\
Given surface $F(x,y,z)=0$:\\
The normal vector gives by the following equation:
\begin{equation}
\boxed{
\nabla F(x,y,z)=\dfrac{\partial F}{\partial x}\mathbf{\hat{i}}+\dfrac{\partial F}{\partial y}\mathbf{\hat{j}}+\dfrac{\partial F}{\partial z}\mathbf{\hat{k}}
}
\end{equation}
substitute $(x_0,y_0,z_0)$ to determine the normal vector at point $(x_0,y_0,z_0)$.
\subsection{Example}
Find the tangent plane of the surface at $(1,0,5)$:
\begin{center}
$xe^{yz}=1$
\end{center}
Solution:\\
Follow thee equation:\\
the tangent vector can be obtained:
\begin{align*}
\nabla F(x,y,z)&=\dfrac{\partial F}{\partial x}\mathbf{\hat{i}}+\dfrac{\partial F}{\partial y}\mathbf{\hat{j}}+\dfrac{\partial F}{\partial z}\mathbf{\hat{k}}\\
&= e^{yz}\mathbf{\hat{i}}+xze^{yz}\mathbf{\hat{j}}+xye^{yz}\mathbf{\hat{k}}\\
&= \mathbf{\hat{i}}+5\mathbf{\hat{j}}+0\mathbf{\hat{k}}
\end{align*}
Therefore, the equation of tangent plane:
\begin{center}
$(x-1)+5(y-0)=0$
\end{center}
\section{Maximum and Minimum Values}
Definition: (For) z=f(x,y) defined in the a region $\mathcal{D}$ in the $xy$-plane)\\
$f(x,y)$ has local maximum (minimum) at (a,b) if:
\begin{center}
$f(x,y)\leq f(a,b)$
\end{center}
for all points $(x,y)$ in the neighbourhood of $(a,b)$.\\
If this relation holds true to all points in the region $\mathcal{D}$, then $f(x,y)$ has an absolute maximum (minimum) at $(a,b)$.
\subsection{Theorem}
If f has a total extremum at $(a,b)$, the first order partial derivatives exist at $(a,b)$, then:
\begin{center}
$f_x(a,b)=f_y(a,b)=0$
\end{center}
\subsection{Critical Point for Extremum}
\begin{itemize}
\item point that satisfy both $f_x=0$ and $f_y=0$.
\item one or both partial derivatives do not exist.
\end{itemize}
local extremums must occur at those point, which one called critical points.
\subsection{Example}
Find all the critical point of this function:
\begin{center}
$f(x,y)=y\sqrt{x}-y^2-x+6y\qquad (x\geq 0)$
\end{center}
Solution:\\
\begin{align*}
f_x &=y\dfrac{1}{2\sqrt{x}}-1\\
f_y &=\sqrt{x}-2y+6
\end{align*}
From those formula, we can see that:\\
$f_x$ doesn't exist at $x=0$,thus ant point lay on the $y$-axis is critical point.\\
Then:
\begin{align*}
f_x &=y\dfrac{1}{2\sqrt{x}}-1=0\\
f_y &=\sqrt{x}-2y+6=0
\end{align*}
this gives another critical point $(4,4)$
\section{Second Derivative Test}
Suppose that the second derivatives of $f(x,y)$ are continuous in a disc centered at $(a,b)$, $f_x(a,b)=f_y(a,b)=0$\\
let:
\begin{center}
$\boxed{\mathcal{D}=
\begin{vmatrix}
f_{xx}(a,b) & f_{xy}(a,b)\\
f_{xy}(a,b) & f_{yy}(a,b)
\end{vmatrix}
}$
\end{center}
Then the second derivative test gives:
\begin{itemize}
\item $\mathcal{D}>0$ $\&$ $f_{xx}(a,b)>0$ $\rightarrow$ Local minimum\\
\item $\mathcal{D}>0$ $\&$ $f_{xx}(a,b)<0$ $\rightarrow$ Local maximum\\
\item $\mathcal{D}<0$ $\rightarrow$ Sudden point \\
\item $\mathcal{D}=0$ $\rightarrow$ Failure
\end{itemize}
\subsection{Absolute Test}
For the absolute extremum, we have not only carry out the second derivative test but also the boundary test.\\
For $y=f(x,y)$ defined in a region $\mathcal{D}$, then $f(x,y)$ has the absolute maximum or absolute minimum at either critical point or some point on the boundary of $\mathcal{D}$

\chapter{Multiple Integral}
\section{Double Integral}
\subsection{Definition}
Recall in two dimensional case:\\
The integral is defined as following:
\begin{center}
$\displaystyle\int_{a}^{b}f(x)=\lim_{\Delta x\to 0}\displaystyle\sum_{i=1}^{n} f(x_i^*)\Delta x_i$
\end{center}
Then, move on to the higher dimension case, say three dimension case:\\
Consider $f(x,y)$ defined in a region $\mathcal{R}$ in the $xy$-plane.
\begin{center}
$\displaystyle\int_{c}^{d}\displaystyle\int_{a}^{b}f(x,y)dxdy =\lim_{\substack{\Delta x\to 0\\ \Delta x\to 0}}\displaystyle\sum_{j=1}^{m}\displaystyle\sum_{i=1}^{n} f(x_i^*,y_j^*)\Delta A_{ij}$
\end{center}
If the limit exists,we called it the double integral for this function within the region $\mathcal{R}$ and denoted by:
\begin{equation}
\boxed{\iint_{\mathcal{R}}f(x,y)dA=\iint_{\mathcal{R}}f(x,y)dxdy}
\end{equation}
In order to evaluate the double integral by the definite integral (For the single variables integral), we introduce by double iterated integral directly from definite integrals.\\
consider $f(x,y)$. if we hold $y$ as a constant, then we may regard $f(x,y)$ as a function of single variables x, then:
\begin{align*}
\displaystyle\int_{a}^{b}f(x,y)dx &= A(y)\\
\displaystyle\int_{c}^{d}A(y)dy &= \displaystyle\int_{c}^{d}\left(\displaystyle\int_{a}^{b}f(x,y)dx\right)dy\\
\displaystyle\int_{c}^{d}f(x,y)dy &= B(x)\\
\displaystyle\int_{a}^{b}B(x)dx &= \displaystyle\int_{a}^{b}\left(\displaystyle\int_{c}^{d}f(x,y)dy\right)dx\\
\end{align*}
Generalized result in two different evaluation process corresponding to the simplification:
\begin{center}
$\boxed{\displaystyle\int_{c}^{d}\left(\displaystyle\int_{h_1(y)}^{h_2(y)}f(x,y)dx\right)dy \quad \& \quad \displaystyle\int_{a}^{b}\left(\displaystyle\int_{g_1(x)}^{g_2(x)}f(x,y)dy\right)dx}$
\end{center}
\subsection{Evaluation in Different Cases}
\begin{itemize}
\item Rectangular Case:
\begin{align*}
\iint_{\mathcal{R}}f(x,y)dA &= \lim_{\substack{\Delta x\to 0\\ \Delta x\to 0}}\displaystyle\sum_{j=1}^{m}\displaystyle\sum_{i=1}^{n} f(x_i^*,y_j^*)\Delta x_{i}\Delta y_{j}\\
&= \lim_{\Delta y\to 0}\displaystyle\sum_{j=1}^{m}\left(\lim_{\Delta x\to 0}\displaystyle\sum_{i=1}^{n} f(x_i^*,y_j^*)\Delta x_{i} \right)\Delta y_{j}\\
&= \lim_{\Delta y\to 0}\displaystyle\sum_{j=1}^{m}\left(\displaystyle\int_{a}^{b}f(x,y)dx \right)\Delta y_{j}\\
&= \displaystyle\int_{c}^{d}\displaystyle\int_{a}^{b}f(x,y)dxdy
\end{align*}
Therefore, in rectangular case, we can convert the double integral to the double iterated integral:
\begin{align*}
\iint_{\mathcal{R}}f(x,y)dA =\quad \cdots \quad=\displaystyle\int_{c}^{d}\displaystyle\int_{a}^{b}f(x,y)dxdy \\
\iint_{\mathcal{R}}f(x,y)dA =\quad \cdots \quad=\displaystyle\int_{a}^{b}\displaystyle\int_{c}^{d}f(x,y)dydx
\end{align*}
These two relationship hold true only for the rectangular region.\\
Furthermore, Fubini's Theorem is based on the rectangular region
$$\mathcal{R}=[a,b]\times [c,d]$$
then the double integral can be evaluated in a simple way:
$$\iint_{\mathcal{R}}f(x)g(y)dA=\int_{a}^{b}f(x)dx\int_{c}^{d}g(y)dy$$
\item $y=h_{i}(x)$ bounded region case:\\
In this case, region $\mathcal{R}={(x,y)|a\leq x\leq b, h_{1}(x)\leq y\leq h_{2}(x)}$
\begin{align*}
\iint_{\mathcal{R}}f(x,y)dA &= \lim_{\substack{\Delta x\to 0\\ \Delta x\to 0}}\displaystyle\sum_{i=1}^{n}\displaystyle\sum_{j=1}^{m} f(x_i^*,y_j^*)\Delta y_{j}\Delta x_{i}\\
&= \lim_{\Delta x\to 0}\displaystyle\sum_{i=1}^{n}\left(\lim_{\Delta y\to 0}\displaystyle\sum_{j=1}^{m} f(x_i^*,y_j^*)\Delta y_{j} \right)\Delta x_{i}\\
&= \lim_{\Delta x\to 0}\displaystyle\sum_{i=1}^{n}\left(\displaystyle\int_{h_{1}(x)}^{h_{2}(x)}f(x,y)dy \right)\Delta x_{i}\\
&= \displaystyle\int_{a}^{b}\left(\displaystyle\int_{h_{1}(x)}^{h_{2}(x)}f(x,y)dy \right)dx
\end{align*}
\item $x=g_{i}(y)$ bounded region case:\\
In this case, region $\mathcal{R}={(x,y)|g_{1}(x)\leq x\leq g_{2}(x), c\leq y\leq d}$
\begin{align*}
\iint_{\mathcal{R}}f(x,y)dA &= \lim_{\substack{\Delta x\to 0\\ \Delta y\to 0}}\displaystyle\sum_{i=1}^{n}\displaystyle\sum_{j=1}^{m} f(x_i^*,y_j^*)\Delta x_{i}\Delta y_{j}\\
&= \displaystyle\int_{c}^{d}\left(\displaystyle\int_{g_{1}(y)}^{g_{2}(y)}f(x,y)dx \right)dy
\end{align*}
\item Two well-defined boundary region case (two point closed loop):\\
In this case, region $\mathcal{R}={(x,y)|g_{1}(x)\leq x\leq g_{2}(x), c\leq y\leq d}$or\\
$\mathcal{R}={(x,y)|a\leq x\leq b, h_{1}(x)\leq y\leq h_{2}(x)}$\\
\begin{align*}
\iint_{\mathcal{R}}f(x,y)dA &= \lim_{\substack{\Delta x\to 0\\ \Delta y\to 0}}\displaystyle\sum_{i=1}^{n}\displaystyle\sum_{j=1}^{m} f(x_i^*,y_j^*)\Delta y_{j}\Delta x_{i}\\
&= \displaystyle\int_{a}^{b}\left(\displaystyle\int_{h_{1}(x)}^{h_{2}(x)}f(x,y)dy \right)dx\\
&= \displaystyle\int_{c}^{d}\left(\displaystyle\int_{g_{1}(y)}^{g_{2}(y)}f(x,y)dx \right)dy
\end{align*}
\end{itemize}
Note:
\begin{itemize}
\item For an irregular regions, we divided it into regular regions.
\begin{align*}
\iint_{\mathcal{R}}f(x,y)dA=\iint_{\mathcal{R}_1}f(x,y)dA+\iint_{\mathcal{R}_2}f(x,y)dA
\end{align*}
\item A double integral (DI) can be evaluated ny a double iterated integral (DII) while DII can be evaluated by some DI
\end{itemize}
\subsection{Properties of DI}
The equations:
\begin{enumerate}
\item $\displaystyle\iint_{\mathcal{R}}dA=$ area of $\mathcal{R}$
\item if $\mathcal{R}=\mathcal{R}_1+\mathcal{R}_2$, then:
$$\iint_{\mathcal{R}}f(x,y)dA=\iint_{\mathcal{R}_1}f(x,y)dA+\iint_{\mathcal{R}_2}f(x,y)dA$$
\item $\displaystyle\iint_{\mathcal{R}}[f(x,y)\pm g(x,y)]dA=\displaystyle\iint_{\mathcal{R}}f(x,y)dA\pm \displaystyle\iint_{\mathcal{R}}g(x,y)dA$
\item $\displaystyle\iint_{\mathcal{R}}cf(x,y)dA=c\displaystyle\iint_{\mathcal{R}}f(x,y)dA$
\end{enumerate}
The inequality:
\begin{enumerate}
\item if $m\leq f(x,y) \leq M$, for all $(x,y)$ in $\mathcal{R}$:
$$mA(\mathcal{R})\leq \iint_{\mathcal{R}}f(x,y)dA \leq MA(\mathcal{R})$$
\end{enumerate}
\section{Double Integral in Polar Coordinates}
Consider the following function:
$$\iint_{\mathcal{R}}f(x,y)dA$$
take an arbitrary point $(x_{ij}^*,y_{ij}^*)$, say $(x_{ij}^*,y_{ij}^*)=(r_{i}\cos \theta_{j},r_{i}\sin \theta_{j})$, then the integral can be rewritten as the following:
\begin{align*}
\iint_{\mathcal{R}}f(x,y)dA &= \lim_{\substack{\Delta r\to 0\\ \Delta \theta\to 0}}\displaystyle\sum_{i=1}^{n}\displaystyle\sum_{j=1}^{m}f(r_{i}\cos \theta_{j},r_{i}\sin \theta_{j})\Delta A_{ij}\\
&=\cdots =f(r_{i}\cos \theta_{j},r_{i}\sin \theta_{j})
\end{align*}
Generalization:
As we can use the Jacobian Matrix to generalize the gradient of the scalar field. The Jacobian can also be thought of as describing the amount of "stretching", "rotating" or "transforming" that a transformation imposes locally. \\
Therefore, in the polar coordinates case:
\begin{align*}
dxdy& =\det\mathcal{J}(r,\theta)drd\theta\\
&=\det
\begin{vmatrix}
\dfrac{\partial x}{\partial r} & \dfrac{\partial x}{\partial \theta} \\
\dfrac{\partial y}{\partial r} & \dfrac{\partial y}{\partial \theta} \\
\end{vmatrix}drd\theta\\
&=\det
\begin{vmatrix}
\cos \theta & -r\sin \theta \\
\sin \theta & r\cos \theta \\
\end{vmatrix}drd\theta\\
&=rdrd\theta
\end{align*}
Thus we obtain:
$$dxdy=rdrd\theta$$
Therefore, after substitute into the definition of the double integral in the polar coordinates:
\begin{itemize}
\item $\displaystyle\iint_{\mathcal{R}}f(x,y)dA=\int_{\alpha_{1}}^{\alpha_{2}}\int_{\rho_{1}}^{\rho_{2}}f(r_{i}\cos \theta_{j},r_{i}\sin \theta_{j})rdrd\theta$
\item $\displaystyle\iint_{\mathcal{R}}f(x,y)dA=\int_{\alpha_{1}}^{\alpha_{2}}\int_{h_{1}(\theta)}^{h_{2}(\theta)}f(r_{i}\cos \theta_{j},r_{i}\sin \theta_{j})rdrd\theta$
\item $\displaystyle\iint_{\mathcal{R}}f(x,y)dA=\int_{\rho_{1}}^{\rho_{2}}\int_{g_{1}(r)}^{g_{1}(r)}f(r_{i}\cos \theta_{j},r_{i}\sin \theta_{j})rd\theta dr$
\end{itemize}
Note that: Usually, if the region involves a circle or part of a circle, we use the polar coordinate to the evaluation, otherwise we use the rectangular coordinates.
\subsection{Example}
Evaluate:
$$\iint_{\mathcal{R}}xdA$$
where the region $\mathcal{R}$ is bounded by $x^2+y^2=2x$\\ \\
Solution: First we find the polar coordinate representation of $\mathcal{R}$:\\
$$r^2=2r\cos \theta$$
then, we start to evaluate the integral:
\begin{align*}
\iint_{\mathcal{R}}xdA &= \int_{-\frac{\pi}{2}}^{\frac{\pi}{2}}\int_{0}^{2\cos \theta}r\cos \theta rdrd\theta \\
&= \int_{-\frac{\pi}{2}}^{\frac{\pi}{2}} \left(\dfrac{\cos \theta}{3}r^3\right)|_{0}^{2\cos \theta}d\theta \\
&= \dfrac{8}{3}\int_{-\frac{\pi}{2}}^{\frac{\pi}{2}}\cos^4\theta d\theta \\
&= \dfrac{8}{3}\int_{-\frac{\pi}{2}}^{\frac{\pi}{2}}\left(\dfrac{1+\cos 2\theta}{2}\right)^2d\theta \\
&= \dfrac{2}{3}\int_{-\frac{\pi}{2}}^{\frac{\pi}{2}}\left(1+2\cos 2\theta +\cos^2 2\theta \right)^2d\theta =\cdots
\end{align*}
\section{Applications of Double Integral}
\subsection{Area}
$$\iint_{\mathcal{R}}dA=A(\mathcal{R})$$
\subsection{Mass of a Plate}
Consider a plate with density varies in the density function $\rho (x,y)$, find the mass of this plate:
\begin{align*}
M &= \lim_{\substack{\Delta x\to 0 \\ \Delta y\to 0}}\sum_{j=1}^{m}\sum_{i=1}^{n}\rho (x_{ij}^*,y_{ij}^*)\Delta A_{ij}\\
&= \iint_{\mathcal{R}}\rho (x,y)dA
\end{align*}
\subsection{Moment of a Plate}
After finding the mass of a plate with density function $\rho (x,y)$, we can further find the moment of a plate:\\
\begin{enumerate}
\item Moment about the $x$-direction
$$\iint_{\mathcal{R}}y\rho (x,y)dA$$
\item Moment about the $x$-direction
$$\iint_{\mathcal{R}}x\rho (x,y)dA$$
\item Center of mass $(x,y)$
$$m\overline{x}=M_y \implies \overline{x}=\dfrac{M_y}{m}=\dfrac{\displaystyle\iint_{\mathcal{R}}x\rho (x,y)dA}{\displaystyle\iint_{\mathcal{R}}\rho (x,y)dA}$$
$$m\overline{y}=M_x \implies \overline{y}=\dfrac{M_x}{m}=\dfrac{\displaystyle\iint_{\mathcal{R}}y\rho (x,y)dA}{\displaystyle\iint_{\mathcal{R}}\rho (x,y)dA}$$
\end{enumerate}
\subsection{Moment of Inertia}
\begin{enumerate}
\item moment of inertia about $x$-axis
$$\iint_{\mathcal{R}}y^2\rho (x,y)dA$$
\item moment of inertia about $y$-axis
$$\iint_{\mathcal{R}}x^2\rho (x,y)dA$$
\item moment of inertia about origin\
$$\iint_{\mathcal{R}}(x^2+y^2)\rho (x,y)dA$$
\subsection{Radii of Gyration}
$$\bar{\bar{x}}^2=\dfrac{\displaystyle\iint_{\mathcal{R}}x^2\rho (x,y)dA}{\displaystyle\iint_{\mathcal{R}}\rho (x,y)dA}$$
$$\bar{\bar{y}}^2=\dfrac{\displaystyle\iint_{\mathcal{R}}y^2\rho (x,y)dA}{\displaystyle\iint_{\mathcal{R}}\rho (x,y)dA}$$
\end{enumerate}
\section{Surface Area}
Given a surface $z=f(x,y)$, the following equation gives the surface area:
\begin{center}
$$\mathcal{S}=\iint_{\mathcal{R}}\sqrt{1+|\nabla f(x,y)|^2}dA$$\\
$$\mathcal{S}=\iint_{\mathcal{R}}\sqrt{1+\dfrac{\partial f}{\partial x}^2+\dfrac{\partial f}{\partial y}^2}dA$$
\end{center}
\section{Triple Integral}
Consider a three dimensional function $f(x,y,z)$ which is defined in a region $\mathcal{D}$ on the space bounded by surface $\mathcal{S}$. We divided $\mathcal{D}$ into many small cube in a similar manner to the double integral, then wecan define the triple integral:
$$\boxed{\iiint_{\mathcal{D}}f(x,y,z)dv=\lim_{\substack{\Delta x\to 0\\ \Delta y\to 0\\ \Delta z\to 0}}\sum_{k=1}^{q}\sum_{j=1}^{m}\sum_{i=1}^{n}f(x_{ijk}^*,y_{ijk}^*,z_{ijk}^*)\Delta x_{i}\Delta y_{j}\Delta z_{k}}$$
\textbf{Generalization:}\\
By considering the boundary condition, the triple integral can be evaluated in three manners:\\
Note: Usually, the three manners is chosen regard to which the projection region is easy to find.
\begin{itemize}
\item $\displaystyle\iiint_{\mathcal{D}}f(x,y,z)dv=\displaystyle\iint_{\mathcal{R}}\left(\int_{\phi_1(x,y)}^{\phi_2(x,y)}f(x,y,z)dz\right)dA$
\item $\displaystyle\iiint_{\mathcal{D}}f(x,y,z)dv=\displaystyle\iint_{\mathcal{R}}\left(\int_{\psi_1(x,z)}^{\psi_2(x,z)}f(x,y,z)dy\right)dA$
\item $\displaystyle\iiint_{\mathcal{D}}f(x,y,z)dv=\displaystyle\iint_{\mathcal{R}}\left(\int_{g_1(y,z)}^{g_2(y,z)}f(x,y,z)dx\right)dA$
\end{itemize}
\section{Triple Integral in Cylindral Spherical Coordinates}
Similar to the tranformation in the double integral, the volume unit is given by the Jacobian matrix determinant.
$$dxdydz=\det \mathcal{J}(r,\theta,z)drd\theta dz=rdrd\theta dz$$
$$dxdydz=\det \mathcal{J}(\rho,\theta,\phi)d\rho d\theta d\phi=\rho ^2\sin \phi d\rho d\theta d\phi$$
Therefore, substitute into the original definition, we can evaluate the triple integral in both cylindral and spherical coordinates:
\begin{align*}
\iiint_{\mathcal{R}}f(x,y,z)dv &=\iiint_{\mathcal{R}}f(r\cos \theta,r\sin \theta,z)rdrd\theta dz \\
&= \iiint_{\mathcal{R}}f(r\sin \phi \cos \theta,r\sin \phi \sin \theta,r\cos \phi)\rho ^2\sin \phi d\rho d\theta d\phi
\end{align*}
\section{Application of Triple Integral}
\subsection{Volume}
$$\iiint_{\mathcal{D}}dv=V(\mathcal{D})$$
\subsubsection{Mass}
In a region $\mathcal{D}$, with the density function $\rho (x,y,z)$, the mass can be given by the following equation:
$$\iiint_{\mathcal{D}}\rho (x,y,z)dv=M(\mathcal{D})$$
\subsection{Mass Center and Moment}
$$\bar{x}=\dfrac{m_{yz}}{M}=\dfrac{\displaystyle\iiint_{\mathcal{D}}x\rho(x,y,z)}{\displaystyle\iiint_{\mathcal{D}}\rho (x,y,z)dv}$$
$$\bar{y}=\dfrac{m_{xz}}{M}=\dfrac{\displaystyle\iiint_{\mathcal{D}}y\rho(x,y,z)}{\displaystyle\iiint_{\mathcal{D}}\rho (x,y,z)dv}$$
$$\bar{z}=\dfrac{m_{xy}}{M}=\dfrac{\displaystyle\iiint_{\mathcal{D}}z\rho(x,y,z)}{\displaystyle\iiint_{\mathcal{D}}\rho (x,y,z)dv}$$
\subsection{Radii of Gyration}
$$\bar{\bar{x}}^2=\dfrac{I_{yz}}{M}=\dfrac{\displaystyle\iiint_{\mathcal{D}}(y^2+z^2)\rho(x,y,z)}{\displaystyle\iiint_{\mathcal{D}}\rho (x,y,z)dv}$$
$$\bar{\bar{y}}^2=\dfrac{I_{xz}}{M}=\dfrac{\displaystyle\iiint_{\mathcal{D}}(x^2+z^2)\rho(x,y,z)}{\displaystyle\iiint_{\mathcal{D}}\rho (x,y,z)dv}$$
$$\bar{\bar{z}}^2=\dfrac{I_{xy}}{M}=\dfrac{\displaystyle\iiint_{\mathcal{D}}(x^2+y^2)\rho(x,y,z)}{\displaystyle\iiint_{\mathcal{D}}\rho (x,y,z)dv}$$
Please refer to the link:\\
\url{http://math.ucsd.edu/~lni/math20e/schedule.html}\\
\url{https://web.math.rochester.edu/people/faculty/edummit/handouts.html}
\chapter{Vector Calculus}
\chapter{Second Order Differential Equation}

\appendix %% CityU never care for you.

\printnotes*

\end{document}
